\section{ماژول a3\_compressor (A3 Compressor)}
\label{sec:module_a3_compressor}

\subsection{معرفی و انگیزه}
این مکانیزم بر مبنای مبانی نظری پیشرفته در ادبیات توجه و فشرده‌سازی \cite{ref_a3compressor_2024} توسعه یافته است.

فشرده‌سازی مدل‌های زبانی از طریق تجزیه مقادیر منفرد (\lr{SVD}) روشی کلاسیک برای کاهش پارامترها و تسریع محاسبات به شمار می‌رود. با این وجود، اعمال مستقیم \lr{SVD} سنتی بر روی ماتریس وزن‌ها بدون در نظر گرفتن ماهیت داده‌های ورودی (\lr{Data-Agnostic}) موجب افت شدید دقت مدل می‌گردد، زیرا برخی از ابعاد ویژگی فعال‌سازی ورودی دارای مقادیر پرت (\lr{Outliers}) و انرژی بالایی هستند که نادیده گرفتن آن‌ها باعث تخریب بازنمایی در لایه‌های عمیق می‌شود.

موتور فشرده‌سازی \lr{A3 Compressor (Activation-Aware SVD with Activation Energy Weighting)} برای حل این معضل توسعه یافته است. در این الگوریتم، به جای تکیه بر اطلاعات خروجی و محاسبات پرهزینه ماتریس اطلاعات فیشر (\lr{Fisher Information Matrix}) که نیازمند گذر پسرو (\lr{Backward Pass}) است، ماتریس وزن‌ها با استفاده از بردار گشتاور دوم خام فعال‌سازی‌های ورودی پیش‌شرط‌سازی (\lr{Preconditioned}) می‌شود. این راهکار به همراه کالیبراسیون مجدد لایه‌های نرمال‌سازی (\lr{LayerNorm Recalibration}) و جستجوی پویای رتبه (\lr{STRS}) امکان کاهش فیزیکی پارامترها بدون فروپاشی بازنمایی مدل را فراهم می‌سازد.

\subsection{مبانی نظری و فرمول‌بندی ریاضی}
فرض کنید $\mathbf{W} \in \mathbb{R}^{d_{\text{out}} \times d_{\text{in}}}$ ماتریس وزن یک لایه خطی متراکم و $\mathbf{X} \in \mathbb{R}^{B \cdot N \times d_{\text{in}}}$ تانسور فعال‌سازی‌های ورودی استخراج‌شده از داده‌های کالیبراسیون باشد. بردار انرژی فعال‌سازی با گشتاور دوم به صورت زیر محاسبه می‌شود:
\begin{equation}
\mathbf{s}_j = \sqrt{\mathbb{E}[x_j^2] + \epsilon} = \sqrt{\frac{1}{M}\sum_{i=1}^M \mathbf{X}_{i, j}^2 + \epsilon}, \quad \forall j \in \{1, \dots, d_{\text{in}}\}
\end{equation}
سپس ماتریس وزن‌ها با قطر این بردار مقیاس‌گذاری می‌شود تا ویژگی‌های پرانرژی تقویت گردند:
\begin{equation}
\tilde{\mathbf{W}} = \mathbf{W} \operatorname{diag}(\mathbf{s})
\end{equation}
بر روی ماتریس مقیاس‌یافته، تجزیه مقادیر منفرد فشرده اجرا می‌شود:
\begin{equation}
\tilde{\mathbf{W}} = \mathbf{U} \mathbf{\Sigma} \mathbf{V}^\top
\end{equation}
فاکتورهای کم‌مرتبه متناظر با رتبه انتخابی $r \ll \min(d_{\text{in}}, d_{\text{out}})$ با توزیع متقارن مقادیر منفرد و اعمال معکوس مقیاس ورودی تشکیل می‌شوند:
\begin{equation}
\mathbf{U}_r = \mathbf{U}_{:, :r} \mathbf{\Sigma}_{:r}^{1/2}, \quad \mathbf{V}_r = \mathbf{\Sigma}_{:r}^{1/2} \mathbf{V}_{:, :r}^\top \operatorname{diag}(\mathbf{s})^{-1}
\end{equation}
در نتیجه بازسازی موثر وزن‌ها به فرم فاکتورگیری‌شده کم‌مرتبه درمی‌آید:
\begin{equation}
\hat{\mathbf{W}} = \mathbf{U}_r \mathbf{V}_r \approx \mathbf{W}
\end{equation}
در گام پایانی، کالیبراسیون مجدد لایه‌های نرمال‌سازی به منظور هماهنگ‌سازی واریانس فعال‌سازی‌های جدید صورت می‌گیرد:
\begin{equation}
\boldsymbol{\gamma}_{\text{new}} = \boldsymbol{\gamma}_{\text{old}} \odot \frac{\boldsymbol{\sigma}_{\text{new}}}{\boldsymbol{\sigma}_{\text{old}}}
\end{equation}

\paragraph{تحلیل طیفی و نامساوی‌های پایداری خطا}
برای بررسی میزان دقت تقریب ماتریس وزن یا نگاشت توجه در ماژول \lr{a3\_compressor}، تفاضل ماتریس فاکتورگیری‌شده از نگاشت متراکم اصلی بر اساس نامساوی هافدینگ-فون‌نیومان ارزیابی می‌گردد. فرض کنید $\mathbf{W} \in \mathbb{R}^{d_{\text{in}} \times d_{\text{out}}}$ ماتریس وزن هدف باشد که توسط تقریب کم‌مرتبه یا پراکنده با رتبه موثر $r$ جایگزین شده است. کران بالای خطای بازسازی بر مبنای نرم فروبنیوس به صورت زیر تصریح می‌شود:
\begin{equation}
\label{eq:spectral_bound_a3_compressor}
\| \mathbf{W} - \mathbf{W}_{r} \|_F = \sqrt{\sum_{i=r+1}^{\min(d_{\text{in}}, d_{\text{out}})} \sigma_i^2(\mathbf{W})} \le \epsilon \| \mathbf{W} \|_F
\end{equation}
که در آن $\sigma_i(\mathbf{W})$ بیانگر $i$-امین مقدار منفرد مرتب‌شده ماتریس به صورت نزولی است. نسبت حفظ انرژی طیفی $\eta_r = \frac{\sum_{i=1}^r \sigma_i^2}{\sum_{i=1}^{d} \sigma_i^2}$ تضمین می‌کند که با انتخاب آستانه $\tau \ge 0.90$، بیش از نود درصد از پراکندگی واریانس داده‌ها حفظ شده و افت سیگنال در انتشار گرادیان رو به عقب به حداقل ممکن کاهش می‌یابد.

\paragraph{تحلیل مجانبی و پیچیدگی اشتراک حافظه سخت‌افزاری}
از منظر پیچیدگی محاسباتی و انتقال داده‌ها در سلسومراتب حافظه (حافظه سراسری \lr{HBM} به حافظه محلی \lr{SRAM})، عملیات ضرب ماتریسی در این مکانیزم نیازمند تخصیص بافر موقت با ابعاد فاکتورگیری‌شده است. برای پردازش دنباله‌ای به طول $N$ با بعد مخفی $D$، بار محاسباتی هسته از فرم درجه دو به رابطه زیر تقلیل می‌یابد:
\begin{equation}
\label{eq:complexity_a3_compressor}
\mathcal{C}_{\text{a3\_compressor}} = \mathcal{O}\left( B \cdot N \cdot D \cdot r + B \cdot H \cdot N \cdot d_k \right) \ll \mathcal{O}\left( B \cdot H \cdot N^2 \cdot d_k \right)
\end{equation}
این بهبود مرتبه پیچیدگی، با کاهش چشمگیر حجم تبادل تانسورها در گذرگاه‌های ارتباطی کارت گرافیک، مانع از اشباع پهنای باند حافظه (پدیده \lr{Memory-Bound}) در فرآیند استنتاج طولانی‌مدت می‌گردد.

\subsection{معماری و پیاده‌سازی}
کلاس موتور فشرده‌سازی \lr{A3Compressor} و لایه خطی فاکتورگیری‌شده \lr{A3FactorizedLinear} در فایل \lr{Codes/framework/compression\_engine/a3\_compressor.py} قرار داشته و با شناسه‌های \lr{a3} و \lr{a3\_compressor} در رجیستری فشرده‌سازی سیستم به ثبت رسیده‌اند. ویژگی‌های برجسته پیاده‌سازی عبارتند از:
\begin{itemize}
    \item \textbf{جایگزینی فیزیکی ماژول‌ها (\lr{True Module Replacement}):} لایه‌های \lr{nn.Linear} و \lr{Conv1D} با ماژول اختصاصی \lr{A3FactorizedLinear} جایگزین می‌شوند که در آن خروجی به شکل $y = (x \mathbf{V}_r^\top)\mathbf{U}_r^\top + b$ محاسبه شده و حجم پارامترها از $d_{\text{in}} \cdot d_{\text{out}}$ به $r \cdot (d_{\text{in}} + d_{\text{out}})$ کاهش می‌یابد.
    \item \textbf{جستجوی تطبیقی رتبه (\lr{STRS Rank Search}):} الگوریتم رتبه‌های نامزد را بررسی کرده و کمترین رتبه‌ای که در آن خطای نسبی فروبنیوس از آستانه مجاز (\lr{max\_error}) کمتر باشد را انتخاب می‌نماید.
    \item \textbf{مسیر امن بازگشت به وزن‌های چگال (\lr{Safety Fallback}):} در صورتی که خطای بازسازی فراتر از آستانه مجاز باشد، ماژول لایه را به ساختار چگال بازگردانده تا از تخریب ساختاری مدل جلوگیری نماید.
    \item \textbf{پشتیبانی از سفیدسازی چولسکی (\lr{Cholesky Whitening}):} تعبیه امکان کواواریانس سفیدسازی کامل مشابه با الگوریتم \lr{SVD-LLM} جهت فشرده‌سازی لایه‌های حساس.
\end{itemize}

\subsection{نتایج تجربی}
ارزیابی بنچ‌مارک سخت‌افزاری و آزمون دقت فشرده‌ساز \lr{a3\_compressor} بر روی مدل‌های گوناگون بر پایه پایگاه اعتبارسنجی \lr{thesis\_final/master\_verified\_results.json} در جدول \ref{tab:benchmark_a3_compressor} منعکس شده است:

\begin{table}[H]
\centering
\small
\caption{نتایج تجربی و اعتبارسنجی بنچ‌مارک فشرده‌ساز \lr{a3\_compressor} روی مدل‌های پایه}
\label{tab:benchmark_a3_compressor}
\begin{tabular}{lccccc}
\toprule
\textbf{مدل پایه} \lr{(Base Model)} & \textbf{دقت} \lr{(Accuracy)} & \textbf{شاخص} \lr{(F1-Score)} & \textbf{حافظه} \lr{(Peak VRAM)} & \textbf{توان} \lr{(Throughput)} & \textbf{وضعیت} \\
\midrule
\lr{DistilBERT} & \lr{77.98\%} & \lr{0.7975} & \lr{773.88 MB} & \lr{26,884.14 tok/s} & \lr{PASS} \\
\lr{RoBERTa} & \lr{85.67\%} & \lr{0.8731} & \lr{389.89 MB} & \lr{8,249.33 tok/s} & \lr{PASS} \\
\lr{GPT-2} & \lr{75.34\%} & \lr{0.7898} & \lr{468.84 MB} & \lr{10,877.23 tok/s} & \lr{PASS} \\
\lr{Qwen} & \lr{51.99\%} & \lr{---} & \lr{1,260.31 MB} & \lr{---} & \lr{PASS} \\
\bottomrule
\end{tabular}
\end{table}

بر اساس نتایج تجربی، ماژول \lr{a3\_compressor} کارایی قابل توجه‌ای در مدل‌های متراکم نشان داده است؛ به طوری که در مدل \lr{RoBERTa} دقت بسیار بالای \lr{85.67\%} همراه با شاخص F1 معادل \lr{0.8731} به ثبت رسیده و در مدل‌های \lr{DistilBERT} و \lr{GPT-2} نیز به ترتیب دقت‌های پایدار \lr{77.98\%} و \lr{75.34\%} گزارش شده است که همگی وضعیت موفقیت‌آمیز \lr{PASS} را احراز کرده‌اند.

\subsection{محدودیت‌ها و تحلیل شکست}
بررسی فرآیند فشرده‌سازی و خطاهای رخ‌داده بیانگر نکات تحلیلی زیر است:
\begin{itemize}
    \item \textbf{افت عملکرد نسبی در مدل \lr{Qwen} (دقت \lr{51.99\%}):} در مدل زبانی \lr{Qwen} دقت برابر با \lr{51.99\%} حاصل شده است. علت اصلی این پدیده، وجود وزن‌های متمرکز با مقادیر پرت بسیار شدید در لایه‌های چندگروهی (\lr{GQA}) مدل‌های مدرن است؛ پیش‌شرط‌سازی ساده با گشتاور دوم در غیاب داده‌های کالیبراسیون غنی و متنوع زبانی، نتوانسته جهت‌گیری تمام کانال‌های پرت را به طور کامل حفظ کند و منجر به تخریب نسبی فاکتورهای کم‌مرتبه در لایه‌های انتهایی دیکودر شده است.
    \item \textbf{حساسیت به فرآیند کالیبراسیون و انحراف لایه‌نرمال:} در صورتی که متن‌های کالیبراسیون ورودی توزیع ناکافی داشته باشند یا پارامتر \lr{recalibrate\_ln} غیرفعال گردد، انحراف واریانس فعال‌سازی در طول لایه‌های عمیق تشدید شده و موجب افزایش خطای فروبنیوس فراتر از حد مجاز (\lr{max\_error}) می‌شود که در آن صورت ماژول ناگزیر به استفاده از مسیر پشتیبان (\lr{fallback}) و بازگردانی لایه به حالت متراکم می‌گردد.
\end{itemize}
